\section{Background}
The limitations of this report are primarily defined by the project requirements from Aalborg University. One of the key requirements was that the solution had to be developed as an embedded system. There is also a need to consider those who will use the technology. In this report, \textit{user} and \textit{wearer} are used interchangeably to describe the person using the hardware, i.e. the wristband. The same is true for the terms \textit{relative} and \textit{guardian}, which are used interchangeably to characterize the person receiving the alerts regarding the person with the wristband. 

Despite the projection that there will be more elderly people in Denmark (see \autoref{fig:pop2070}), a study by Ældresagen states that the population will not necessarily consist of more pensioners. Part of the reasoning behind this statement has to do with the rising retirement age. According to projections by Ældre Sagen, there will be 58,000 more elderly people above retirement age in 2040 compared to the number of individuals above retirement age in 2024.

Another part of the limitations is set by how the users view on the use of technology in healthcare. In a study from 2025 \cite{noauthor_undersogelse_nodate}, the population's attitude towards municipal use of welfare technology for elderly care was investigated.

The study was conducted by Epinion using a combination of web-based and telephone interviews. It is representative with respect to gender, age, region, and educational level, and a total of 1,603 responses were collected from individuals aged 18 and above.

Three requirements for welfare technology were found: 

\begin{enumerate}
\item \textbf{You need to choose yourself. }
This is in regard to a primarily positive attitude towards the use of welfare technology; however, with a stance that the individual should have the possibility to select or deselect welfare technology.

\item \textbf{Fear of loneliness. }
This refers to a worry about whether or not the municipal use of welfare technology causes loneliness among elderly people, due to the municipal staff's fewer visits. 

\item \textbf{The responsibility of municipalities in the use of welfare technology. }This reflects the general opinion that the elderly person should not have to pay for welfare themselves if the municipality stops offering the welfare to the people in need. This opinion stretches to the attitude that the responsibility of the municipality also includes paying, installing, operations and data security in terms of the use of welfare technology \cite{noauthor_undersogelse_nodate}.
\end{enumerate}

In the article, \textit{Undersøgelse: Befolkningens holdning til velfærdsteknologi} \cite{noauthor_undersogelse_nodate}, a set of limitations to technology were identified. What is also important to note is the experiences of older people with technology, thereby their abilities and lack thereof. A study by Epinion has delved into what experiences and challenges the elderly demographic experiences when using technology. In the study 35\% agree that the challenges they experience with technology are, to some extent, great or very high. Of these, 10\% say the challenges they face are high or very high; this corresponds to approximately half a million people. 

The study also takes a look at the population's view of how digitalisation impacts them. To this, 40\% of people aged 75 or older feel forced to incorporate digital solutions in their everyday lives without wanting to do so. In terms of younger demographics, 29\% of 55-64-year-olds, 26\% of 35-54-year-olds and 28\% of 18-34-year-olds said the same thing. This leads to a need for digital help, guidance and assistance. 71\% of the population experience needing digital help, whereas 41\% of the population feel that they need this guidance occasionally or more often.

Besides the technological challenges mentioned by this study, 86\% of the population agrees or strongly agrees that being able to talk to a human being is also very important in addition to digital self-service solutions. 

% As many as 41\% of the population feel they need help, advice, or guidance from others to use digital devices or solutions occasionally or more often
% 71\% of the population may feel they need digital assistance

\subsection{Problem statement}
That is why, in this project, an embedded fall detection system is designed and evaluated with the aim of supporting elderly individuals living independently through timely detection of falls and notification of relatives. Therefore, the following problem statement was formulated:


\textbf{How can embedded systems and fall detection algorithms be used to support elderly citizens living independently by detecting falls and ensuring timely notification of relatives?}

To accompany this statement, the following sub questions have been formulated.
\begin{itemize}
    \item Which fall detection algorithms are suitable for implementation in low-power wearable devices?
    \item Which sensors are best suited for fall detection in an embedded system?
    \item How can a fall detection product be designed to ensure adoption and usability among elderly users?
\end{itemize}

\subsection{The Sustainable Development Goals}
In addition to the previously mentioned limitations from the project description regarding the product, there was also a project requirement by the university, that the solution should help one of the UN's Sustainable Development Goals. The goal which the solution aids is goal number 3, which is about health and well-being for everyone of all ages. Specifically, the solution helps sub-point number 8, which focuses on universal health coverage, which also includes access to quality vital health care assistance \cite{martin_health_nodate}. 

As previously mentioned, it is estimated that in the future there will be a very unequal distribution of young people and elderly people in Denmark. Considering that the older one gets, more expenses arise for the healthcare system. This adds a financial pressure on the healthcare system in the form of less money per senior citizen. The solution in this report aims to ensure that more elderly people have the option of staying at home rather than in a nursing home, despite sometimes needing help. In addition to this, it is the relatives who are the first in line of support for the elderly person. This results in saving money on healthcare costs and working hours of healthcare staff. This leads to more resources for the expected increase in the number of elderly people. The fact that elderly people can live at home also strengthens their retirement as they retain a great deal of independence and are in familiar surroundings, which can increase their quality of life. 

 